% Compile with pdfLaTeX. One column; searchable text; no external assets.
% Tailor to each vacancy using only skills and experience you can substantiate.
\documentclass[11pt,letterpaper]{article}
\usepackage[utf8]{inputenc}
\usepackage[T1]{fontenc}
\usepackage{mathptmx}
\usepackage[left=0.5in,right=0.5in,top=0.48in,bottom=0.48in]{geometry}
\usepackage{enumitem}
\usepackage[hidelinks]{hyperref}
\usepackage{microtype}
\input{glyphtounicode}
\pdfgentounicode=1
\hypersetup{pdftitle={Roberto Guagnelli Garcia - Mechanical Engineering CV},pdfauthor={Roberto Guagnelli Garcia}}
\pagestyle{empty}
\setlength{\parindent}{0pt}
\setlength{\parskip}{0pt}
\hyphenpenalty=10000
\exhyphenpenalty=10000
\raggedright
\setlist[itemize]{leftmargin=13pt,labelsep=5pt,topsep=2pt,itemsep=1.5pt,parsep=0pt,partopsep=0pt}
\newcommand{\cvsection}[1]{\vspace{8pt}{\large\bfseries #1}\par\vspace{1pt}\hrule\vspace{4pt}}
\newcommand{\entry}[2]{#1\hfill #2\par}
\begin{document}
{\Large\bfseries ROBERTO GUAGNELLI GARCÍA}\par
\vspace{2pt}
{\small +52 55 4976 5729 $|$ \href{mailto:robguagnelli@gmail.com}{robguagnelli@gmail.com} $|$ Juriquilla, Querétaro, Mexico}\par
{\small \href{https://www.linkedin.com/in/roberto-guagnelli-garcia/}{linkedin.com/in/roberto-guagnelli-garcia} $|$ Portfolio: \href{https://robertoguagnelli.github.io/}{robertoguagnelli.github.io}}\par
\vspace{3pt}
\textbf{Mechanical Engineering Student $|$ Mechanical Design, CAD \& Aircraft Integration}

\cvsection{EDUCATION}
\entry{\textbf{Tecnológico de Monterrey}, Campus Estado de México}{Expected Jul 2027}
B.S. in Mechanical Engineering; Concentration in Aerospace Engineering\par
\textbf{Academic average: 92/100 (3.7/4.0 GPA)}

\cvsection{WORK EXPERIENCE}
\entry{\textbf{GE Aerospace} $|$ Designer Intern, Product Definition}{Aug 2026 -- Present}
\textit{RISE Open Fan Program}
\begin{itemize}
\item Develop electrical harness routing and mounting bracket CAD models in Siemens NX; review interferences and apply design for manufacturability (DFM) and maintainability criteria to component integration.
\item Prepare engineering drawings defining component geometry and design requirements for the RISE Open Fan program.
\end{itemize}
\vspace{3pt}
\entry{\textbf{GE Aerospace} $|$ Lean Summer Experience Intern}{Jun 2026 -- Jul 2026}
\textit{Integrated Performance \& Operability (IP\&O)}
\begin{itemize}
\item Implemented visual planning, dashboards and agreed review practices with subsection managers and program leads; reduced average weekly meeting duration by 20\%, as timed by meeting leaders.
\item Defined and tracked key performance indicators (KPIs) and supported cloud database migration to improve performance visibility across management and engineering teams.
\end{itemize}

\cvsection{ENGINEERING PROJECTS \& LEADERSHIP}
\entry{\textbf{Borregos CEM AeroDesign} $|$ Tecnológico de Monterrey}{2024 -- Present}
\entry{\textit{Engineering Coordinator / Design Captain}}{Aug 2024 -- Present}
\begin{itemize}
\item Lead a five-person design team and coordinate 40+ members across design, structures, electronics and R\&D; developed Wingull (2025) and Latías (2026) from conceptual design through flight testing.
\item Personally modeled 300+ CAD parts per aircraft within 400+ component assemblies; evaluated mass, center of gravity, stability, materials and motor/propeller combinations to guide design decisions.
\item Performed steady-state and transient computational fluid dynamics (CFD) studies of airfoils and vortex formation to assess aerodynamic performance and inform stability evaluation; used drag coefficients in takeoff models to estimate payload capability at different density altitudes.
\item Designed removable joints for Wingull's 4.2 m wing to enable breakdown into components under 1.2 m; supported an 18 kg distributed-mass wing test with measured tip deflection within 8\% of predictions.
\item Redesigned a 3D-printed PETG bottle retainer for Latías after insertion-force and fracture issues; revised geometry completed 50 insertion/removal cycles without observed cracks or permanent deformation.
\item Integrated Latías' twin-engine configuration, which flew with 6 kg payload at 7 kg empty mass; Wingull flew with 4 kg payload at 6 kg empty mass.
\end{itemize}
\vspace{3pt}
\entry{\textit{Manufacturing Coordinator $|$ Hyperion}}{2024}
\begin{itemize}
\item Coordinated eight members in aircraft fabrication, assembly and quality control; designed and manufactured a welded-aluminum nose landing gear with a spring, travel stop and fork to accommodate limited propeller clearance.
\item Supported six competition flights carrying 3 kg payload at 4 kg empty mass; team placed 2nd overall and 3rd in design report at SAE Aero Design México 2024.
\end{itemize}
\cvsection{TECHNICAL SKILLS}
{\small
\textbf{CAD \& Analysis:} Siemens NX, SolidWorks, Fusion 360, ANSYS Fluent, XFLR5; steady-state and transient CFD.\par
\textbf{Programming \& Data:} Python, MATLAB, Excel, Power BI.\par
\textbf{Design \& Testing:} Engineering drawings, interference checks, DFM, maintainability, propulsion selection, structural and impact testing.\par
\textbf{Manufacturing:} 3D printing, carbon-fiber fabrication, laser cutting, aluminum welding, tube and sheet bending.\par
\textbf{Certification:} Certified SOLIDWORKS Professional (CSWP), perfect exam score.\par
\textbf{Languages:} Spanish (native), English (C2; Cambridge and IELTS), German (A2).
}
\end{document}
